\section{Step 30 Pause between the first update and the next read}\label{crc-s30}

Stop on the second visit to \texttt{file\_read}, before it starts another read. With the three-byte fixture, \texttt{s5} holds raw \texttt{77b79c2d} and the descriptor is still open. Then continue to the return boundary and check the finalized result and descriptor cleanup. Saved registers are restored on return, so this working state must be inspected inside the function.

Use \texttt{s30-1}, \texttt{s30-2}. These CLI fixture observations pause on the second visit to \texttt{file\_read}, then continue to the return boundary.

\subsection{Expected observation}
A positive short read advances the raw state but does not finalize or close the file.

\Needspace{11\baselineskip}
\subsection{Check your understanding}
Why inspect \texttt{s5} before restoration?

\begin{itemize}
  \item \textbf{A} It is the caller result register.
  \item \textbf{B} The restoring exit reinstates the caller saved value.
  \item \textbf{C} The state is stored in the filename.
\end{itemize}
